% GENERATED FILE. Edit canonical lessons, then run npm run generate:books.
% canonical-source-hash: fnv1a64:48c9a9bf48bffd50
% canonical-lessons: LA-C483-stagnum, LA-C483-gurges, LA-C483-vorago, LA-C483-mano, LA-C483-inundo

\chapter{Pool, Whirlpool, Chasm, To trickle, To flood}
\label{ch:la-pool-whirlpool-chasm-to-trickle-to-flood}
\hlchaptermodality{483}

\begin{chapteropening}
\textbf{By the end of this chapter:} \emph{I can talk about still and dangerous water: a pool, a whirlpool, a chasm, water trickling and a river flooding.}

Complete the last lesson of chapter 483: I can talk about still and dangerous water: a pool, a whirlpool, a chasm, water trickling and a river flooding.
\end{chapteropening}

\section[stagnum, stagnī]{stagnum, stagnī — a pool, a pond; standing water}

\label{lesson:LA-C483-stagnum}



\begin{quote}
Before the new one: say the Latin for to flow, then the Latin for to swim across.
\end{quote}

\subsection*{You'll want to know: stagnum, stagnī}
\textbf{stagnum, stagnī} — \textquotedblleft{}a pool, a pond; standing water\textquotedblright{}.

\begin{quote}
In silvā stagnum quiētum est, et rānae cantant.
\end{quote}

\begin{quote}
\emph{In the wood there is a quiet pool, and frogs sing.}
\end{quote}

Second declension, neuter: \emph{stagnum, stagnī n.} Still water, not running. English \emph{stagnant} comes from the verb \emph{stagnāre}.

\begin{cousinweb}
The origin of \emph{stagnum} is not settled: Lewis \& Short compare Greek \emph{tenagos} 'shallow water', but EtymDB gives no parent. English \emph{stagnant} comes from \emph{stagnāre}, built on it.
\end{cousinweb}

\begin{grammarlens}[title={What you've built}]
One more word for telling, explaining and describing at B1.
\end{grammarlens}

\subsection*{Your turn}
\begin{itemize}
\raggedright
  \item \emph{Say it:} \emph{stagnum, stagnī}
  \item \emph{Say it:} \emph{stagnum, stagnī}, once more
  \item \emph{Say it:} a sentence of your own with \emph{stagnum, stagnī}
  \item \emph{Recall:} say the Latin for to flow, then the Latin for to swim across, then say \emph{stagnum, stagnī} again
\end{itemize}

\begin{tcolorbox}[breakable,colback=teal!4,colframe=teal!35!black,title={Before you move on}]
What does stagnum, stagnī mean? (\textquotedblleft{}A pool, a pond; standing water\textquotedblright{}.) Say it once more.
\end{tcolorbox}
\section[gurges, gurgitis]{gurges, gurgitis — a whirlpool, an abyss of water; a flood}

\label{lesson:LA-C483-gurges}



\begin{quote}
Before the new one: say the Latin for to swim across, then the Latin for a pool.
\end{quote}

\subsection*{You'll want to know: gurges, gurgitis}
\textbf{gurges, gurgitis} — \textquotedblleft{}a whirlpool, an abyss of water; a flood\textquotedblright{}.

\begin{quote}
Nauta in gurgitem cadit, sed amīcī eum servant.
\end{quote}

\begin{quote}
\emph{A sailor falls into a whirlpool, but his friends save him.}
\end{quote}

Third declension, masculine: \emph{gurges, gurgitis m.} Also a glutton, one who swallows everything. English \emph{regurgitate} comes from it.

\begin{cousinweb}
\emph{gurges} is inherited and Lewis \& Short link it with \emph{gula} 'throat'; EtymDB lists Greek \emph{gorgyra} 'underground drain' as a cousin. English \emph{regurgitate} comes from a word built on it.
\end{cousinweb}

\begin{grammarlens}[title={What you've built}]
One more word for telling, explaining and describing at B1.
\end{grammarlens}

\subsection*{Your turn}
\begin{itemize}
\raggedright
  \item \emph{Say it:} \emph{gurges, gurgitis}
  \item \emph{Say it:} \emph{gurges, gurgitis}, once more
  \item \emph{Say it:} a sentence of your own with \emph{gurges, gurgitis}
  \item \emph{Recall:} say the Latin for to swim across, then the Latin for a pool, then say \emph{gurges, gurgitis} again
\end{itemize}

\begin{tcolorbox}[breakable,colback=teal!4,colframe=teal!35!black,title={Before you move on}]
What does gurges, gurgitis mean? (\textquotedblleft{}A whirlpool, an abyss of water; a flood\textquotedblright{}.) Say it once more.
\end{tcolorbox}
\section[vorāgō, vorāginis]{vorāgō, vorāginis — a chasm, an abyss; a whirlpool}

\label{lesson:LA-C483-vorago}



\begin{quote}
Before the new one: say the Latin for a pool, then the Latin for a whirlpool.
\end{quote}

\subsection*{You'll want to know: vorāgō, vorāginis}
\textbf{vorāgō, vorāginis} — \textquotedblleft{}a chasm, an abyss; a whirlpool\textquotedblright{}.

\begin{quote}
In mediō forō terra sē aperit, et vorāgō magna fit.
\end{quote}

\begin{quote}
\emph{In the middle of the forum the ground opens, and a great chasm forms.}
\end{quote}

Third declension, feminine: \emph{vorāgō, vorāginis f.}, from \emph{vorō} 'to swallow'.

\begin{cousinweb}
\emph{vorāgō} is built on \emph{vorō} 'to swallow, to devour' (Lewis \& Short): a hole that swallows things up. English \emph{voracious} comes from \emph{vorō}.
\end{cousinweb}

\begin{grammarlens}[title={What you've built}]
One more word for telling, explaining and describing at B1.
\end{grammarlens}

\subsection*{Your turn}
\begin{itemize}
\raggedright
  \item \emph{Say it:} \emph{vorāgō, vorāginis}
  \item \emph{Say it:} \emph{vorāgō, vorāginis}, once more
  \item \emph{Say it:} a sentence of your own with \emph{vorāgō, vorāginis}
  \item \emph{Recall:} say the Latin for a pool, then the Latin for a whirlpool, then say \emph{vorāgō, vorāginis} again
\end{itemize}

\begin{tcolorbox}[breakable,colback=teal!4,colframe=teal!35!black,title={Before you move on}]
What does vorāgō, vorāginis mean? (\textquotedblleft{}A chasm, an abyss; a whirlpool\textquotedblright{}.) Say it once more.
\end{tcolorbox}
\section[mānō, mānāre]{mānō, mānāre — to trickle, to flow, to drip}

\label{lesson:LA-C483-mano}



\begin{quote}
Before the new one: say the Latin for a whirlpool, then the Latin for a chasm.
\end{quote}

\subsection*{You'll want to know: mānō, mānāre}
\textbf{mānō, mānāre} — \textquotedblleft{}to trickle, to flow, to drip\textquotedblright{}.

\begin{quote}
Aqua frīgida ē saxō mānat.
\end{quote}

\begin{quote}
\emph{Cold water trickles out of the rock.}
\end{quote}

\emph{mānō, mānāre, mānāvī, mānātum}. Also of news: to spread. English \emph{emanate} comes from \emph{ēmānāre}.

\begin{cousinweb}
The origin of \emph{mānō} is not settled: Lewis \& Short connect it, as probable, with \emph{madeō} 'to be wet', but EtymDB gives no parent. English \emph{emanate} comes from the compound \emph{ēmānāre}.
\end{cousinweb}

\begin{grammarlens}[title={What you've built}]
One more word for telling, explaining and describing at B1.
\end{grammarlens}

\subsection*{Your turn}
\begin{itemize}
\raggedright
  \item \emph{Say it:} \emph{mānō, mānāre}
  \item \emph{Say it:} \emph{mānō, mānāre}, once more
  \item \emph{Say it:} a sentence of your own with \emph{mānō, mānāre}
  \item \emph{Recall:} say the Latin for a whirlpool, then the Latin for a chasm, then say \emph{mānō, mānāre} again
\end{itemize}

\begin{tcolorbox}[breakable,colback=teal!4,colframe=teal!35!black,title={Before you move on}]
What does mānō, mānāre mean? (\textquotedblleft{}To trickle, to flow, to drip\textquotedblright{}.) Say it once more.
\end{tcolorbox}
\section[inundō, inundāre]{inundō, inundāre — to flood, to overflow}

\label{lesson:LA-C483-inundo}



\begin{quote}
Before the new one: say the Latin for a chasm, then the Latin for to trickle.
\end{quote}

\subsection*{You'll want to know: inundō, inundāre}
\textbf{inundō, inundāre} — \textquotedblleft{}to flood, to overflow\textquotedblright{}.

\begin{quote}
Vēre flūmen agrōs inundat.
\end{quote}

\begin{quote}
\emph{In spring the river floods the fields.}
\end{quote}

\emph{inundō, inundāre, inundāvī, inundātum}, from \emph{in-} and \emph{unda} 'wave'. English \emph{inundate} comes from it.

\begin{cousinweb}
\emph{inundō} is built from \emph{in-} and \emph{undō} 'to rise in waves', from \emph{unda} 'wave' (EtymDB; Lewis \& Short). English \emph{inundate} comes from it.
\end{cousinweb}

\begin{grammarlens}[title={What you've built}]
One more word for telling, explaining and describing at B1.
\end{grammarlens}

\subsection*{Your turn}
\begin{itemize}
\raggedright
  \item \emph{Say it:} \emph{inundō, inundāre}
  \item \emph{Say it:} \emph{inundō, inundāre}, once more
  \item \emph{Say it:} a sentence of your own with \emph{inundō, inundāre}
  \item \emph{Recall:} say the Latin for a chasm, then the Latin for to trickle, then say \emph{inundō, inundāre} again
\end{itemize}

\begin{tcolorbox}[breakable,colback=teal!4,colframe=teal!35!black,title={Before you move on}]
What does inundō, inundāre mean? (\textquotedblleft{}To flood, to overflow\textquotedblright{}.) Say it once more.
\end{tcolorbox}
